% Part: first-order-logic
% Chapter: sequent-calculus
% Section: soundness-identity

\documentclass[../../../include/open-logic-section]{subfiles}

\begin{document}

\olfileid[id]{fol}{seq}{sid}

\olsection{Kesahihan dengan \usetoken{S}{identity}}

\begin{prop}
$\Log{LK}$ dengan sekuen awal dan aturan identitas bersifat sahih.
\end{prop}

\begin{proof}
Sekuen awal berbentuk ${} \Sequent \eq[t][t]$ valid, sebab untuk setiap
!!{structure}~$\Struct M$, berlaku $\Sat{M}{\eq[t][t]}$. (Perhatikan
bahwa kita mengandaikan suku $t$ tertutup, yakni tidak memuat variabel,
sehingga penugasan variabel tidak relevan.)

Andaikan inferensi terakhir dalam !!a{derivation} adalah $=$. Maka
premisnya ialah $\eq[t_1][t_2], \Gamma \Sequent \Delta, !A(t_1)$ dan
kesimpulannya ialah $\eq[t_1][t_2], \Gamma \Sequent \Delta, !A(t_2)$.
Perhatikan !!a{structure}~$\Struct M$. Kita perlu menunjukkan bahwa
kesimpulannya valid; artinya, jika $\Sat{M}{\eq[t_1][t_2]}$ dan
$\Sat{M}{\Gamma}$, maka entah $\Sat{M}{!C}$ untuk suatu
$!C \in \Delta$ atau $\Sat{M}{!A(t_2)}$.

Menurut hipotesis induksi, premisnya valid. Artinya, jika
$\Sat{M}{\eq[t_1][t_2]}$ dan $\Sat{M}{\Gamma}$, maka entah (a) untuk
suatu $!C \in \Delta$, $\Sat{M}{!C}$, atau (b)
$\Sat{M}{!A(t_1)}$. Dalam kasus (a), pembuktian selesai. Perhatikan kasus
(b). Misalkan $s$ adalah penugasan variabel dengan
$s(x) = \Value{t_1}{M}$. Menurut
\olref[syn][ass]{prop:sentence-sat-true}, $\Sat{M}{!A(t_1)}[s]$. Karena
$\varAssign{s}{s}{x}$, menurut \olref[syn][ext]{prop:ext-formulas},
$\Sat{M}{!A(x)}[s]$. Karena $\Sat{M}{\eq[t_1][t_2]}$, kita mempunyai
$\Value{t_1}{M} = \Value{t_2}{M}$, sehingga
$s(x) = \Value{t_2}{M}$. Dengan menerapkan
\olref[syn][ext]{prop:ext-formulas} sekali lagi, kita juga memperoleh
$\Sat{M}{!A(t_2)}[s]$. Menurut
\olref[syn][ass]{prop:sentence-sat-true}, $\Sat{M}{!A(t_2)}$.
\end{proof}

\end{document}
